% Generated draft: requires mathematical and source review.
\documentclass[11pt]{article}
\usepackage[margin=25mm]{geometry}
\usepackage{fontspec}
\setmainfont{Latin Modern Roman}
\usepackage{amsmath,amssymb,mathtools,unicode-math}
\setmathfont{Latin Modern Math}
\usepackage{enumitem,needspace,xurl,hyperref}
\setlength{\parindent}{0pt}
\setlength{\parskip}{5pt}
\setlength{\emergencystretch}{3em}
\newcommand{\N}{\mathbb{N}}
\newcommand{\Z}{\mathbb{Z}}
\newcommand{\Q}{\mathbb{Q}}
\newcommand{\R}{\mathbb{R}}
\newcommand{\C}{\mathbb{C}}
\newcommand{\F}{\mathbb{F}}
\newcommand{\A}{\mathbb{A}}
\newcommand{\PP}{\mathbb{P}}
\newcommand{\E}{\mathbb{E}}
\newcommand{\eps}{\varepsilon}
\newcommand{\dd}{\,\mathrm{d}}
\begin{document}
\section*{UnsolvedMath ID: 30002157}
\textbf{Source problem code:} catalogue entry
\subsection{Definitions and mathematical statement}
Let \(p\) be a prime, let \(F\) be a finite extension of \(\Q_p\), and let \(E/F\) be an unramified quadratic extension. Write \(\mathcal O_F\) for the ring of integers of \(F\) and \(a\mapsto\bar a\) for the nontrivial automorphism of \(E/F\). The source asks for an identity in every dimension \(n\); it does not explicitly specify whether the dimension range includes \(n=0\).

The symbols \(\mathcal N_n\) and \(\mathcal N_{n+1}\) denote the formal moduli spaces intended in the source's arithmetic fundamental lemma formulation. Their moduli problems, signatures, polarization conditions, framing objects, and base formal schemes are not supplied. If \(i:\mathcal N_n\hookrightarrow\mathcal N_{n+1}\) is the intended embedding, the stated diagonal cycle is its graph
\[
Y=\{(x,i(x)):x\in\mathcal N_n\}\subseteq\mathcal N_n\times\mathcal N_{n+1}.
\]
The notation \(U(n+1)\) requires a choice of a nondegenerate Hermitian \(E\)-space \((V,h)\) of dimension \(n+1\). Its group of \(F\)-rational isometries is
\[
U(V)(F)=\{g\in\operatorname{GL}_E(V):h(gv,gw)=h(v,w)\text{ for all }v,w\in V\}.
\]
The source does not specify this Hermitian space or the action of the relevant group on \(\mathcal N_{n+1}\).

A customary interpretation of the symmetric-space notation appearing on the analytic side is
\[
S_{n+1}(F)=\{g\in\operatorname{GL}_{n+1}(E):g\bar g=I_{n+1}\},\qquad
S_{n+1}(\mathcal O_F)=S_{n+1}(F)\cap\operatorname{GL}_{n+1}(\mathcal O_E),
\]
where conjugation is entrywise and \(\mathcal O_E\) is the ring of integers of \(E\). This convention is not explicitly defined in the supplied record and requires confirmation. Write \(\mathbf 1_A\) for the characteristic function of a set \(A\).

\textbf{Proposed identity, as recorded.} For every dimension \(n\) in the intended range and every pair of matching elements
\[
\delta\in U(n+1),\qquad \gamma\in\operatorname{GL}_{n+1}(E),
\]
where matching is described in the source as being under \(\operatorname{GL}_n(E)\), is it true that
\[
\bigl\langle Y,(1\times\delta)^*Y\bigr\rangle
=
\operatorname{orb}^{\prime}\!\left(\gamma,\mathbf 1_{S_{n+1}(\mathcal O_F)}\right)?
\]
Here the left side is the intended arithmetic intersection number, and the right side is the intended derivative orbital integral. The displayed equality retains the source's notation and normalization without alteration.

\textbf{Unresolved qualifications.} The record does not define matching, specify the intended regularity conditions on the elements, or explain whether \(\gamma\) must belong to \(S_{n+1}(F)\). It also does not define the intersection multiplicity, its finiteness conditions, the orbital-integral deformation parameter, the differentiation point, Haar measures, transfer factors, signs, or logarithmic normalization. Consequently, this is a source-faithful schematic conjectural identity; the supplied information does not determine a fully specified numerical equality. These missing choices must be recovered before a precise arithmetic fundamental lemma statement can be accepted.
\subsection{Short English statement}
For every intended dimension, the problem asks whether the arithmetic intersection of the diagonal cycle with its unitary translate equals a derivative orbital integral for matching elements over an unramified quadratic extension of \(p\)-adic fields. The source omits essential geometric data, admissibility conditions, matching definitions, and normalization conventions, so the displayed identity requires human mathematical review. The supplied label is partially solved; its dated literature assessment has not been independently verified.
\subsection{Sources}
\begin{itemize}
\item Local UnsolvedMath registry, supplied complete record: permanent ID \(30002157\), title \emph{Arithmetic Fundamental Lemma in All Dimensions}, exact source problem code \texttt{catalogue entry}, collection \emph{Oberwolfach Reports Open Problems}. No individual proposer is supplied; the recorded proposed year is 2012.
\item Supplied report citation: \emph{Arithmetic Geometry (2013)}, Oberwolfach Reports, DOI \texttt{10.4171/owr/2012/38}, pp.\ 2335--2388. \url{https://doi.org/10.4171/owr/2012/38}. The background describes these reports as documenting mathematical workshops and their open-problem discussions. The citation's year 2013 and the record's year 2012 are retained without reconciliation.
\item Supplied literature source: Zhang, \emph{Weil representation and AFL}. \url{https://annals.math.princeton.edu/2021/193-3/p05}.
\item Supplied literature source: Mihatsch--Zhang, \emph{AFL over a general p-adic field}. \url{https://arxiv.org/abs/2104.02779}.
\item Supplied literature source: \emph{Oberwolfach report}. \url{https://doi.org/10.4171/owr/2012/38}.
\item Literature assessment dated exactly \texttt{2026-08-21}: supplied dated curation, not independently verified. It labels the problem partially solved and reports proofs in important minuscule, dimension, and \(p\)-adic cases, while describing the refined arithmetic Gan--Gross--Prasad formula and the full general arithmetic fundamental lemma as unresolved. The refined formula is background context, not a separately formulated target here. The source expressly requests independent verification of time-sensitive status and cited claims.
\item Supplied statement status: \texttt{corrected\_verified}. The accompanying verification field says that the official report states the all-dimensional arithmetic fundamental lemma after recording proved low-dimensional and special cases. This is a supplied source-verification assertion; the report and linked literature have not been independently consulted for this draft.
\end{itemize}
\end{document}
